% ============================================================================
\chapter{Srodni radovi}
\label{ch:srodni}
% ============================================================================

U ovom poglavlju daje se pregled dosada\v{s}njih istra\v{z}ivanja koja su bliska
temi rada. Cilj je smjestiti vlastiti rad u \v{s}iri kontekst i jasno pokazati
\v{c}ime se on razlikuje od postoje\'ce literature. Najprije se razmatraju radovi
o mjerenju performansi operativnih sistema za rad u realnom vremenu, zatim radovi
o zamjeni konteksta i prekidima na arhitekturi RISC-V, potom radovi o mekim
procesorima na programabilnoj logici, te na kraju radovi o odnosu simulacije i
stvarnog hardvera. Na osnovu tog pregleda izdvaja se pozicija ovog rada.

\section{Mjerenje performansi operativnih sistema za rad u realnom vremenu}
\label{sec:srodni-rtos}

Mjerenje vremena zamjene konteksta uobi\v{c}ajena je metrika pri pore\dj{}enju
operativnih sistema za rad u realnom vremenu. Pore\dj{}enja sistema FreeRTOS i
Zephyr pokazuju da se ovakve veli\v{c}ine izra\v{z}avaju u broju takta, pri \v{c}emu
se za FreeRTOS navodi red veli\v{c}ine od oko dvije stotine dvadeset tri takta na
jezgrama tipa ARM Cortex-M \citep{ulsolutions2022rtos}. Va\v{z}no je naglasiti da
takve vrijednosti mjere u\v{z}u operaciju zamjene konteksta na cjevovodnom
procesoru, dok mjerenja u ovom radu obuhvataju i dodatni tro\v{s}ak rada sa
mehanizmom sinhronizacije na vi\v{s}eciklusnom mekom procesoru, pa se brojevi ne
mogu direktno pore\dj{}ivati bez te ograde.

\section{Zamjena konteksta i prekidi na arhitekturi RISC-V}
\label{sec:srodni-riscv}

Poseban pravac istra\v{z}ivanja bavi se ubrzanjem zamjene konteksta i obrade
prekida na arhitekturi RISC-V. U radu o pro\v{s}irenju CV32RT predla\v{z}u se
mikroarhitekturna pro\v{s}irenja jezgre radi br\v{z}eg ulaska u prekid i br\v{z}e
zamjene konteksta za mikrokontrolere zasnovane na RISC-V \citep{balas2023cv32rt}.
Sli\v{c}no, postoji istra\v{z}ivanje koje zajedni\v{c}ki istra\v{z}uje
mikroarhitekturu procesora i pro\v{s}irenja samog operativnog sistema FreeRTOS sa
ciljem smanjenja latencije zamjene konteksta \citep{scheck2026coexploration}.
Ovi radovi pokazuju da je smanjenje cijene zamjene konteksta aktuelna tema, ali
se oni uglavnom oslanjaju na izmjene hardvera ili jezgre, dok se ovaj rad bavi
mjerenjem nepromijenjenog sistema na postoje\'coj mekoj jezgri.

\section{Meki procesori RISC-V na programabilnoj logici}
\label{sec:srodni-softcore}

Meki procesori na programabilnoj logici predmet su brojnih evaluacija. Procesor
PicoRV32, kori\v{s}ten i u ovom radu, ocijenjen je kao izrazito \v{s}tedljiv po
pitanju resursa i pogodan za ure\dj{}aje sa ograni\v{c}enim resursima
\citep{jahnke2023picorv32}. Takve evaluacije obi\v{c}no se fokusiraju na
zauze\'ce logi\v{c}kih elemenata i propusnost, dok je rad operativnog sistema na
takvoj jezgri i njegovo uklapanje zajedno sa video izlazom rje\dj{}e obra\dj{}ivano.
Upravo to \v{c}ini doprinos ovog rada, jer se uz mjerenje razmatra i prakti\v{c}no
ograni\v{c}enje blokovske memorije pri sintezi.

\section{Odnos simulacije i stvarnog hardvera}
\label{sec:srodni-qemu}

Emulator QEMU \v{s}iroko se koristi za pokretanje ugradbenog softvera, ali je on
funkcijski ta\v{c}an, \v{s}to zna\v{c}i da vjerno opona\v{s}a pona\v{s}anje
instrukcija, ali ne modeluje mikroarhitekturu takovno ta\v{c}no \citep{qemu}.
Zbog toga se vrijednosti takta dobijene u takvom okru\v{z}enju ne mogu bez ograde
tuma\v{c}iti kao vrijednosti na stvarnom \v{c}ipu. Ova osobina direktno opravdava
opreznost pri pore\dj{}enju rezultata iz simulatora \citep{thesisqemu} sa
rezultatima izmjerenim na hardveru, kako je u\v{c}injeno u poglavlju
\ref{ch:mjerenje}.

\section{Pozicija ovog rada}
\label{sec:srodni-pozicija}

Iz pregleda literature mo\v{z}e se uo\v{c}iti obrazac. Mjerenja performansi
operativnih sistema naj\v{c}e\v{s}\'ce se rade na cjevovodnim mikrokontrolerima
tipa ARM Cortex-M, dok se na arhitekturi RISC-V te\v{z}i\v{s}te stavlja na
hardversko ubrzanje zamjene konteksta. Evaluacije mekih procesora fokusiraju se
na resurse, a rje\dj{}e na rad operativnog sistema. Ovaj rad popunjava prazninu
time \v{s}to mjeri nepromijenjeni FreeRTOS na minimalnoj otvorenoj mekoj RISC-V
jezgri na jeftinoj plo\v{c}i, upore\dj{}uje rezultat sa simulatorom i dokumentuje
prakti\v{c}no ograni\v{c}enje sinteze vezano za blokovsku memoriju. Time se
doprinos rada jasno razlikuje od navedenih istra\v{z}ivanja.

\section{Zaklju\v{c}ak poglavlja}
\label{sec:srodni-zakljucak}

U poglavlju je dat pregled srodnih istra\v{z}ivanja iz \v{c}etiri bliske oblasti
i izdvojena je pozicija ovog rada. Pokazano je da postoje\'ca literatura uglavnom
mjeri operativne sisteme na cjevovodnim jezgrama ili predla\v{z}e hardverska
ubrzanja, dok mjerenje nepromijenjenog sistema na minimalnoj mekoj RISC-V jezgri,
zajedno sa pore\dj{}enjem sa simulatorom, ostaje slabije pokriveno. Time je
opravdan pristup koji se razra\dj{}uje u narednim poglavljima.
